\section{Reconstructing the law with every marginal retained}
\label{p13:transfer:section}

At a precision finer than $N^{-1}$, removing the exceptional windows is
too expensive. We retain their exact marginals instead. This also keeps
the repeated occurrence of a single Bernoulli variable inside the
reference law, where it can be evaluated exactly. The remaining task is
to control interactions between distinct windows and to sum the expansion
without applying a fixed-order estimate at an order growing with $N$.

We use the full admissible set $\mathcal I_N=\{x\in I_N:p_x>0\}$,
the exact marginals $p_x=\E J_x$, the centered coefficients
$c(I)=\E\prod_{x\in I}(J_x-p_x)$, and the relative nullity defined in
Section~\ref{p13:exact:section}. Throughout this section we abbreviate
$\rho_{\mathrm{rel}}(I)$ to $\rho(I)$; this convention uses the effective
local ranks and does not add the local rank defects again. Put
\[
 A_2=\sum_{|e|=2}|c(e)|,\qquad
 A_3=\sum_{\substack{|T|=3\\T\text{ connected}}}|c(T)|,
 \qquad p_* =\max_{x\in\mathcal I_N}p_x.
\]
The marginal and pair estimates give
$p_*\leq2p\ll_C N^{-1}$, $\mu=\sum_xp_x=O_C(1)$, and
$A_2\ll_{C,\varepsilon}N^{-2/3+\varepsilon}$.
The connection between the arithmetic estimates and the distribution is
the following bound:
\begin{equation}
 \|\Law(Z_N)-Q_{\mu_N,\alpha_N}\|_1
 \ll_{C,\varepsilon}N^\varepsilon A_3+N^{-4/3+\varepsilon}.
 \label{p13:transfer:announced-bound}
\end{equation}
We prove it in Theorem~\ref{p13:transfer:reconstruction}. The two terms
have different origins: the first records a genuinely connected
three-window interaction, while the second permits two pair interactions.

\paragraph{Three errors in the reconstruction.}
For $e=\{x,y\}$ write $S_e=S_x+S_y$.
The retained class \(\mathscr D_2\) consists of direct families and
families \(e\cup U\) with a dependent pair \(e\) for which
\(S_e\) and the spaces \((S_z)_{z\in U}\) form a direct sum.
The dependent pair is then unique. Write
\(T_{\mathscr D_2}=\sum_{I\in\mathscr D_2}q(I)\Delta^{*|I|}\).
With the covariance-corrected Bernoulli measure \(B_2\) defined in
\eqref{p13:transfer:B2}, the proof separates
\[
 \begin{aligned}
 \|\Law(Z_N)-Q_{\mu_N,\alpha_N}\|_1
 &\le \|\Law(Z_N)-T_{\mathscr D_2}\|_1
       +\|T_{\mathscr D_2}-B_2\|_1\\
 &\qquad +\|B_2-Q_{\mu_N,\alpha_N}\|_1.
 \end{aligned}
\]
The first term is the stopping boundary, with the unexpanded suffix
kept in a positive conditional measure. The second consists of the
independent and pair deficits in
\eqref{p13:transfer:retained-difference}; the covariance signs are
retained before taking norms. In particular, an incompatible pair
is allowed in \(\mathscr D_2\): its joint probability is zero,
but its covariance is still present in \(B_2\).
The last term costs \(O(N^{-5/3+\varepsilon})\) by
Lemma~\ref{p13:transfer:conversion}; that rate concerns the full
measure \(B_2\), including every covariance.
Theorem~\ref{p13:transfer:finite} assembles the three estimates,
using fixed orders and the separately uniform tail of
Proposition~\ref{p13:transfer:uniform-tail} for the boundary.

\subsection{The exact Bernoulli reference and its conversion}

For any finite family of Bernoulli variables, without an independence
assumption, define $b_i=\delta_0+p_i\Delta$ and
\begin{equation}
 B_2=\mathop{*}_{i\in\mathcal I_N}b_i
       +\sum_{|e|=2}c(e)\Delta^{*2}*
                     \mathop{*}_{i\notin e}b_i.
 \label{p13:transfer:B2}
\end{equation}
This is a signed measure. Its first term is the law of a sum of
independent Bernoulli variables with the exact individual marginals;
its second term restores every covariance once.

\begin{lemma}[Conversion with heterogeneous marginals]
\label{p13:transfer:conversion}
Let $s_j=\sum_i p_i^j$, $C_2=\sum_{|e|=2}c(e)$, and
$\alpha=2C_2-s_2$. If $\mu\leq M$ and $p_*\leq1/4$, then
\begin{equation}
 \|B_2-Q_{\mu,\alpha}\|_1
 \leq C_M\{s_3+s_2^2+A_2(s_2+p_*)\}.
 \label{p13:transfer:conversion-bound}
\end{equation}
For the present model this is $O_{C,\varepsilon}(N^{-5/3+\varepsilon})$.
\end{lemma}

\begin{proof}
In the convolution algebra of finite signed measures, the series for
$\log_*(\delta_0+p_i\Delta)$ converges absolutely, since
$\|p_i\Delta\|_1=2p_i\leq1/2$. Summing it gives
\[
 \log_*\left(\mathop{*}_ib_i\right)
   =\mu\Delta-\frac{s_2}{2}\Delta^{*2}+E,
 \qquad \|E\|_1\leq\sum_{j\geq3}\frac{2^js_j}{j}
                   \leq\frac{16}{3}s_3.
\]
Exponentiation and $\exp_*(\mu\Delta)=\Pi_\mu$ imply
\[
 \left\|\mathop{*}_ib_i-\Pi_\mu
              +\frac{s_2}{2}\Pi_\mu*\Delta^{*2}\right\|_1
 \leq C_M(s_3+s_2^2).
\]
The product omitting a pair $e=\{i,j\}$ differs from
$\Pi_{\mu-p_i-p_j}$ by $O_M(s_2)$ in norm, and
\[
 \|\Pi_{\mu-p_i-p_j}-\Pi_\mu\|_1\leq2(p_i+p_j).
\]
Replacing each such product in \eqref{p13:transfer:B2} by $\Pi_\mu$
therefore costs $O_M(A_2(s_2+p_*))$. The exact identity
$\Var Z-\E Z=2C_2-s_2$ proves the lemma.
Here $s_2\leq p_*\mu=O_C(N^{-1})$ and
$s_3\leq p_*^2\mu=O_C(N^{-2})$, giving the stated application.
\end{proof}

\paragraph{Algebraic form of the conversion.}
The estimate in Lemma~\ref{p13:transfer:conversion} also holds for
arbitrary real pair weights $c(e)$ in \eqref{p13:transfer:B2}, with
$C_2=\sum_e c(e)$, $A_2=\sum_e|c(e)|$ and $\alpha=2C_2-s_2$.
The mean and marginal bounds are unchanged. The two approximations
in its proof and the triangle inequality give this algebraic form;
only the interpretation $\alpha=\Var Z-\E Z$ uses actual covariances.


\subsection{Arithmetic estimates used in the assembly}

The proof below uses three estimates proved in
Section~\ref{p2:technical-transfer} of the companion. They separate
fixed geometry from the growing number of windows.
For each fixed $r\ge4$, the total connected activity is
\[
 \sum_{\substack{|I|=r\\I\text{ connected}}}|c(I)|
 \ll_{r,C,\varepsilon}N^\varepsilon A_3+N^{-4/3+\varepsilon},
\]
and its incompatible part is $O(N^{-3/2+\varepsilon})$.
The independent deficit of all pair extensions is
$O(N^{-3/2+\varepsilon})$ at every cardinality.
Finally, the connected activity beyond the fixed size $256$,
up to the logarithmic boundary cutoff and with compatible nullity
at most $2L$, is $O(N^{-4/3+\varepsilon})$; the corresponding
primitive tail has the same bound without a cardinality cutoff.
These are Proposition~\ref{p13:transfer:fixed-activities},
Proposition~\ref{p13:transfer:pair-deficit} and
Proposition~\ref{p13:transfer:uniform-tail}, respectively.

The last estimate has a separate proof. It uses small-kernel counts,
moderate congruence certificates and a graph with small incoming and
outgoing degrees. It does not let the order in the first estimate
grow with $N$. The choice $256$ is sufficient for these estimates;
no optimality is attached to that cutoff.

\paragraph{Dependencies of the reconstruction.}
The table separates the arithmetic inputs by the quantity they control.
The fixed-order estimates and the uniform tails enter different parts
of the stopped expansion; the lower bound uses the independently proved
moment estimate.
\begin{center}
\small
\begin{tabularx}{\linewidth}{@{}>{\raggedright\arraybackslash}p{0.26\linewidth}>{\raggedright\arraybackslash}X>{\raggedright\arraybackslash}p{0.24\linewidth}@{}}
\toprule
Quantity & Main inputs in the companion & Use in the article\\
\midrule
Complete pair coefficient
& Sections~\ref{p1:arith:appendix}--\ref{p1:arith:supplement}
& Theorem~\ref{p1:arith:pairs}\\[5pt]
Rank of a nonaligned pair
& \ref{p13:rig:short-basis}--\ref{p13:rig:nonaligned};
  norm estimate in \ref{p13:rig:two-witnesses}
& Theorem~\ref{p13:third:main}\\[5pt]
Absolute triple activity
& Section~\ref{p13:activity:section};
  \ref{p13:third:rational-rank}--\ref{p13:third:lower}
& Corollary~\ref{p13:third:activity}\\[5pt]
Connected activity at fixed orders
& \ref{p13:transfer:dense-targets},
  \ref{p13:transfer:primitive},
  \ref{p13:transfer:fixed-deficit},
  \ref{p13:transfer:fixed-activities}
& Theorem~\ref{p13:transfer:reconstruction}\\[5pt]
Pair deficit at every cardinality
& \ref{p13:transfer:independent-deficit}--\ref{p13:transfer:pair-deficit}
& Equation~\eqref{p13:transfer:retained-error}\\[5pt]
Uniform tails
& Section~\ref{p14:cert:section};
  \ref{p13:transfer:independent-deficit},
  \ref{p13:transfer:independent-tail}--\ref{p13:transfer:uniform-tail}
& Theorems~\ref{p13:transfer:finite}--\ref{p13:transfer:reconstruction}\\[5pt]
Quadratic lower bound
& Theorem~\ref{p1:fixed:moments}; incompatible part of
  Proposition~\ref{p13:transfer:fixed-activities}
& \ref{p13:lower:quartic}--\ref{p13:lower:optimal}\\
\bottomrule
\end{tabularx}
\end{center}

\subsection{The stopped expansion and its finite input}

Let $\mathscr D_2$ consist of the direct families and the families
$e\cup U$ for which $e$ is any dependent pair and $S_e,(S_z)_{z\in U}$
are direct. Affine incompatibility of $e$ is allowed. This class is
hereditary; every non-direct member contains exactly one dependent pair.
For a fixed ordering of $\mathcal I_N$, its boundary contains the first
excluded prefixes. A boundary family has relative nullity at most $2L$:
its retained predecessor has at most $L$, and adding one block adds
at most $L$. In particular every subfamily of a boundary family has
compatible nullity at most $2L$.

\begin{theorem}[Finite transfer to the complete law]
\label{p13:transfer:finite}
Fix $1<\beta\leq4/3$. Suppose that, for each of the fixed integers
$3\leq r\leq256$ and every $\varepsilon>0$,
\begin{equation}
 \sum_{\substack{|I|=r,\ I\text{ connected}\\\rho_*(I)\leq2L}}
                   |c(I)|\ll_{r,C,\varepsilon}N^{-\beta+\varepsilon}.
 \label{p13:transfer:finite-hypothesis}
\end{equation}
Then
\begin{equation}
 \|\Law(Z_N)-Q_{\mu_N,\alpha_N}\|_1
                  \ll_{C,\varepsilon}N^{-\beta+\varepsilon}.
 \label{p13:transfer:finite-conclusion}
\end{equation}
Only these $254$ fixed orders are assumed in
\eqref{p13:transfer:finite-hypothesis}; the larger orders are supplied
by Proposition~\ref{p13:transfer:uniform-tail}.
\end{theorem}

\begin{proof}
Fix a weight $w>0$. The hypothesis at sizes at most $256$ and
\eqref{p13:transfer:connected-tail} give, for $K_N=o(L)$,
\begin{equation}
 T:=\sum_{\substack{I\text{ connected},\ 3\leq|I|\leq K_N\\
                             \rho_*(I)\leq2L}}
                  w^{|I|}|c(I)|\ll N^{-\beta+\eta}.
 \label{p13:transfer:connected-range}
\end{equation}
The finitely many factors $w^r$ at $r\leq256$ only change constants.
For a primitive $P$, $\rho(P)\leq L$ and
$\prod_{x\in P}p_x\leq|c(P)|$, because
$c(P)=q(P)-\prod_Pp_x$ and $q(P)$ is either zero or at least twice
the independent product. Together with
\eqref{p13:transfer:primitive-tail}, this gives
\begin{equation}
 \sum_{\substack{P\text{ primitive}\\|P|\geq3}}
                     w^{|P|}|c(P)|\ll N^{-\beta+\eta}
 \label{p13:transfer:all-primitives}
\end{equation}
over all cardinalities.

Put $T_{\mathscr D_2}=\sum_{I\in\mathscr D_2}q(I)\Delta^{*|I|}$.
The finite expansion of $B_2$ yields the identity
\begin{equation}
 T_{\mathscr D_2}-B_2
 =-\sum_{I\notin\mathscr D_2}\Delta^{*|I|}
       \left[\prod_{i\in I}p_i+
          \sum_{\substack{e\subset I\\|e|=2}}
                   c(e)\prod_{i\in I\setminus e}p_i\right].
 \label{p13:transfer:retained-difference}
\end{equation}
For the independent term, choose a primitive subfamily of a discarded
family. If it has at least three members, use
\eqref{p13:transfer:all-primitives}; if it is a pair $e$, the complement
must be a non-direct extension of $S_e$, and
$\prod_{i\in e}p_i\leq|c(e)|$.
The pair terms in \eqref{p13:transfer:retained-difference} are bounded
by the same all-cardinality pair deficit. Free complements have total
weight at most $e^{2\mu}$. Consequently
\begin{equation}
 \|T_{\mathscr D_2}-B_2\|_1\ll N^{-\beta+\eta}.
 \label{p13:transfer:retained-error}
\end{equation}

Take $K_N=(5+o(1))\log N/\log\log N$.
The exact stopping identity from Section~\ref{p1:stop:section} gives
\[
 \|\Law(Z_N)-T_{\mathscr D_2}\|_1
       \leq\sum_{I\in\partial\mathscr D_2}2^{|I|}q(I).
\]
For $|I|\leq K_N$, expand $q(I)$ in centered coefficients. The terms
of size zero and two are the independent and pair deficits already
bounded. Every remaining nonzero centered term has either a connected
component of size at least three or at least two pair components.
Canonical factorization, \eqref{p13:transfer:connected-range}, and
the free independent complements bound their sum by a constant times
\[
 e^{4A_2+T}-1-4A_2=O(T+A_2^2).
\]
Here $T=o(1)$ and $A_2^2\ll N^{-4/3+\eta}$; the centered coefficient
of a singleton is zero. This is where the linear contribution of one
pair is subtracted before taking an absolute bound.

On the larger boundary, $q(I)\leq2^{2L}\prod_Ip_x$, so
\[
 \sum_{\substack{I\in\partial\mathscr D_2\\|I|>K_N}}2^{|I|}q(I)
 \ll_C N^2\sum_{r>K_N}\frac{(2\mu)^r}{r!}=N^{-3+o(1)}.
\]
Combining these bounds, \eqref{p13:transfer:retained-error}, and
Lemma~\ref{p13:transfer:conversion} proves
\eqref{p13:transfer:finite-conclusion}. All internal margins are chosen
after the final $\varepsilon$; none is made to vary with $N$.
\end{proof}

\begin{theorem}[Reconstruction from triple activity]
\label{p13:transfer:reconstruction}
With the exact mean and second factorial cumulant,
\begin{equation}
 \|\Law(Z_N)-Q_{\mu_N,\alpha_N}\|_1
       \ll_{C,\varepsilon}N^\varepsilon A_3+N^{-4/3+\varepsilon}.
 \label{p13:transfer:reconstruction-bound}
\end{equation}
In particular the arithmetic triple estimate gives the upper bound
$O_{C,\varepsilon}(N^{-4/3+\varepsilon})$ for the full model.
\end{theorem}

\begin{proof}
In the proof of Theorem~\ref{p13:transfer:finite}, replace the finite
activity input by Proposition~\ref{p13:transfer:fixed-activities}.
At sizes $3,\ldots,256$ its total cost is
$N^{\varepsilon'}A_3+N^{-4/3+\varepsilon'}$; the larger sizes have the
second cost by Proposition~\ref{p13:transfer:uniform-tail}.
The primitive deficits, pair deficits, and boundary expansion obey
the same bounds. In the component generator,
\[
 T\ll N^{\varepsilon'}A_3+N^{-4/3+\varepsilon'},\qquad
 e^{4A_2+T}-1-4A_2=O(T+A_2^2).
\]
The established arithmetic estimate makes $T=o(1)$, so this passage
loses no new fixed power of $N$. The large boundary costs $N^{-3+o(1)}$
and conversion costs $N^{-5/3+\varepsilon'}$.
To make the choices explicit, denote the requested final margin by
$\varepsilon_{\mathrm{out}}=\varepsilon$ and choose
$0<\tau<\min(\varepsilon_{\mathrm{out}}/4,1/960)$.
Invoke each completed fixed-order estimate with sufficiently small fixed
input margins that its total displayed loss is at most $\tau$.
For example, use the pair estimate with loss $\tau/2$ before squaring
$A_2$. Only the fixed orders $3,\ldots,256$ occur in this choice;
the larger orders use the uniform certificate bounds, not constants
from a fixed-order theorem. With $\varepsilon'=\tau$ the preceding
argument proves \eqref{p13:transfer:reconstruction-bound}, since
$\tau<\varepsilon_{\mathrm{out}}$.

Finally invoke Corollary~\ref{p13:third:activity} with loss $\tau$.
The triple term is then at most $N^{-41/30+2\tau}$, and
\[
 \frac{41}{30}-2\tau-\frac43=\frac1{30}-2\tau>0.
\]
The conversion and large-boundary exponents exceed $4/3$ by $1/3$
and $5/3$, respectively. These comparisons give the announced upper
bound. All accuracy choices, source counts and fixed activity orders used here
are fixed before $N$; the growing cardinality cutoff $K_N$ is controlled
separately.
\end{proof}

\paragraph{Power margins.}
The strict power savings used in the last stages are recorded below.
Each margin compares complete contributions, after coefficient choices,
propagation and rank factors have been included.
\begin{center}
\small
\begin{tabularx}{\textwidth}{@{}Xcl@{}}
\toprule
Comparison & Exact margin & Location\\
\midrule
Nonaligned target above a rational pair versus dominant triples
& $11/8-41/30=1/120$
& \eqref{p13:third:unaligned-target-cost}\\[4pt]
Incompatible triples versus the third factorial cumulant
& $3/2-41/30=2/15$
& \eqref{p13:activity:kappa-comparison}\\[4pt]
Third activity versus the quadratic remainder
& $41/30-4/3=1/30$
& \eqref{p13:transfer:reconstruction-bound}\\[4pt]
Low-nullity uniform tail versus the quadratic remainder
& $27/20-4/3=1/60$
& \eqref{p13:transfer:low-rank-tail}\\[4pt]
Quartic signed error versus $\alpha_N^2$
& $3/2-4/3=1/6$
& \eqref{p13:lower:positive-quartic}\\
\bottomrule
\end{tabularx}
\end{center}
The smallest displayed margin is $1/120$. In contrast, two pair
activities have size $A_2^2\ll N^{-4/3+\eta}$ and are retained at
the target scale. They have no strict power saving to spend.
For a prescribed final accuracy $\varepsilon_0>0$, one may choose
$0<\tau<\min(\varepsilon_0/4,1/960)$ and invoke completed auxiliary
estimates with total loss at most $\tau$ each. Their own input margins
are chosen smaller as needed; in a product, divide the allotted loss
between the factors. Thus the same symbol $\varepsilon$ in successive
statements never requires the same numerical value at every stage.
Rank thresholds and fixed geometric degrees have separate roles.
The fixed connected-activity estimates in reconstruction are used only
through order $256$; larger cardinalities have the uniform certificate estimate and a factorial
boundary bound.

